
\documentclass{article}
\usepackage{colortbl}
\usepackage{makecell}
\usepackage{multirow}
\usepackage{supertabular}

\begin{document}

\newcounter{utterance}

\centering \large Interaction Transcript for game `bbh', experiment `hyperbaton', episode 155 with Qwen3.6{-}27B.
\vspace{24pt}

{ \footnotesize  \setcounter{utterance}{1}
\setlength{\tabcolsep}{0pt}
\begin{supertabular}{c@{$\;$}|p{.15\linewidth}@{}p{.15\linewidth}p{.15\linewidth}p{.15\linewidth}p{.15\linewidth}p{.15\linewidth}}
    \# & \multicolumn{2}{c}{Player} && \multicolumn{2}{c}{Game Master} \\
    \hline

    \theutterance \stepcounter{utterance}  
    & & & \multicolumn{4}{p{0.6\linewidth}}{
        \cellcolor[rgb]{0.9,0.9,0.9}{
            \makecell[{{p{\linewidth}}}]{
                \texttt{\tiny{[P1$\langle$GM]}}
                \texttt{Order adjectives correctly in English sentences.} \\
\\ 
\texttt{Q: Which sentence has the correct adjective order:} \\
\texttt{Options:} \\
\texttt{(A) rubber terrible ship} \\
\texttt{(B) terrible rubber ship} \\
\texttt{A: (B)} \\
\\ 
\texttt{Q: Which sentence has the correct adjective order:} \\
\texttt{Options:} \\
\texttt{(A) repulsive small Brazilian exercise ship} \\
\texttt{(B) Brazilian repulsive exercise small ship} \\
\texttt{A: (A)} \\
\\ 
\texttt{Q: Which sentence has the correct adjective order:} \\
\texttt{Options:} \\
\texttt{(A) blue gold wonderful square shoe} \\
\texttt{(B) wonderful square blue gold shoe} \\
\texttt{A: (B)} \\
\\ 
\texttt{Q: Which sentence has the correct adjective order:} \\
\texttt{Options:} \\
\texttt{(A) square cardboard old{-}fashioned Japanese exercise white motorcycle} \\
\texttt{(B) old{-}fashioned square white Japanese cardboard exercise motorcycle} \\
\texttt{A:} \\
            }
        }
    }
     \\ \\

    \theutterance \stepcounter{utterance}  
    & \multicolumn{4}{p{0.6\linewidth}}{
        \cellcolor[rgb]{0.561,0.737,0.561}{
            \makecell[{{p{\linewidth}}}]{
                \texttt{\tiny{[P1$\rangle$GM]}}
                \texttt{To determine the correct adjective order, we follow the standard English sequence: **Opinion, Size, Age, Shape, Color, Origin, Material, Purpose, Noun**.} \\
\\ 
\texttt{Let's analyze the adjectives in the options:} \\
\texttt{1.  **Opinion**: old{-}fashioned (subjective evaluation)} \\
\texttt{2.  **Size**: (none)} \\
\texttt{3.  **Age**: (none, "old{-}fashioned" is opinion here, not age like "old")} \\
\texttt{4.  **Shape**: square} \\
\texttt{5.  **Color**: white} \\
\texttt{6.  **Origin**: Japanese} \\
\texttt{7.  **Material**: cardboard} \\
\texttt{8.  **Purpose**: exercise} \\
\texttt{9.  **Noun**: motorcycle} \\
\\ 
\texttt{**Correct Order:**} \\
\texttt{Opinion (old{-}fashioned) {-}> Shape (square) {-}> Color (white) {-}> Origin (Japanese) {-}> Material (cardboard) {-}> Purpose (exercise) {-}> Noun (motorcycle).} \\
\\ 
\texttt{Result: **old{-}fashioned square white Japanese cardboard exercise motorcycle**} \\
\\ 
\texttt{Let's check the options:} \\
\texttt{(A) square cardboard old{-}fashioned Japanese exercise white motorcycle {-}> Incorrect order (Shape before Opinion, Material before Opinion, Color at the end).} \\
\texttt{(B) old{-}fashioned square white Japanese cardboard exercise motorcycle {-}> Matches the correct order.} \\
\\ 
\texttt{A: (B)} \\
            }
        }
    }
    & & \\ \\

    \theutterance \stepcounter{utterance}  
    & & & \multicolumn{2}{p{0.3\linewidth}}{
        \cellcolor[rgb]{0.9,0.9,0.9}{
            \makecell[{{p{\linewidth}}}]{
                \texttt{\tiny{[GM$|$GM]}}
                \texttt{Target: (B)} \\
            }
        }
    }
    & & \\ \\

    \theutterance \stepcounter{utterance}  
    & & & \multicolumn{2}{p{0.3\linewidth}}{
        \cellcolor[rgb]{0.9,0.9,0.9}{
            \makecell[{{p{\linewidth}}}]{
                \texttt{\tiny{[GM$|$GM]}}
                \texttt{game\_result = LOSE} \\
            }
        }
    }
    & & \\ \\

\end{supertabular}
}

\end{document}
